% ============================================================
% CHAPTER 2: EXECUTIVE SUMMARY
% MCDA 5585 / 5586 Major Project Report
% Bhavik Kantilal Bhagat | A00494758 | Saint Mary's University
% ============================================================

\chapter{Executive Summary}
\label{chap:executive_summary}

This report documents my five-month Mitacs Business Strategy Internship (BSI)
at Citco Technology Management (CTM), Halifax, Nova Scotia, conducted between
\textbf{March~15 and August~31, 2026} under the project title
\textit{``Intelligent Automation and Reliability Engineering for Financial
Services''} (Mitacs reference \textbf{IT49539}).

% ────────────────────────────────────────────────────────────
\section{The Problem}
% ────────────────────────────────────────────────────────────

Before my internship began, Citco's IA-SRE team operated without a
centralized observability platform. Monitoring for 119 Lambda functions,
22 ECS clusters, 83 EC2 instances, and critical Kafka-backed event pipelines
across five AWS environments relied on an entirely reactive model: business
teams in Dublin, London, and Halifax would notice anomalies in fund processing
outputs and raise tickets. No CloudWatch alarms existed on any Meridian
service. No centralized dashboard surfaced platform health. When the
Meridian event pipeline stalled on April~17, 2026 --- backing up 9~million
Kafka messages over 7.5 hours --- the SRE team learned about it from a
Dublin business user, not from any monitoring system.

Citco also uses \textbf{Dynatrace} for BPM automation observability, but
its enterprise licensing costs scale with infrastructure size. My industry
supervisor, Kishor Deotale, framed the core research question for the Mitacs
engagement: \textit{can AWS-native services achieve Dynatrace-equivalent
observability at a fraction of the cost?}

% ────────────────────────────────────────────────────────────
\section{The Approach}
% ────────────────────────────────────────────────────────────

During this internship, I delivered a \textbf{proof-of-concept (PoC)} AWS-native
observability stack and feasibility assessment, organized across three
parallel work streams:

\begin{enumerate}
    \item \textbf{SRE Observability Platform (IISS-487)} --- A centralized
          Amazon Managed Grafana (AMG) dashboard with 14 rows, 110+ panels,
          and 38 CloudWatch alarms covering all major AWS service categories,
          deployed via parameterized CloudFormation Infrastructure-as-Code.

    \item \textbf{IA Monitoring Platform (IISS-482--486)} --- Per-platform
          CloudWatch dashboards and a registry-based SQS monitoring pipeline,
          establishing the architectural patterns and resource inventory that
          the Grafana platform builds on.

    \item \textbf{IA Resources Inventory API (IISS-507)} --- A Flask REST
          API for programmatic discovery of all cloud resources across five
          budget code environments, with 562 automated tests at 99\% coverage.
\end{enumerate}

Four dedicated RPA production support rotations (Weeks 8, 12, 16, 20) ran
in parallel, where I handled live incidents including the CAIS Deadline credential
hotfix (IISS-706) and the MESO month-end stuck tasks (IISS-741, Critical).

% ────────────────────────────────────────────────────────────
\section{Key Results}
% ────────────────────────────────────────────────────────────

\begin{table}[ht]
\centering
\caption{Executive Summary --- Key Quantitative Outcomes}
\label{tab:exec_outcomes}
\begin{tabular}{p{6.5cm}p{2.5cm}p{3.0cm}}
\toprule
\textbf{Metric} & \textbf{Before} & \textbf{After} \\
\midrule
Incident detection time (MSK lag) & 2--3 hours & $<$30 seconds \\
CloudWatch alarms on IA platform & 0 & 38 \\
Centralized dashboard panels & 0 & 110+ \\
Container Insights (Meridian ECS) & 1 of 6 clusters & 5 of 6 clusters \\
Resource inventory automation & Manual/Confluence & 119 Lambda, 22 ECS, 83 EC2 (API) \\
Inventory API test coverage & N/A & 562 tests, 99\% \\
AWS-native vs Dynatrace parity & N/A & 60\% (80\% with App Signals) \\
Monthly observability cost (AWS) & \$49 (CW only) & \$77 (CW + Grafana + X-Ray) \\
\bottomrule
\end{tabular}
\end{table}

The most operationally significant outcome of my work is the reduction in Kafka consumer
lag detection from a multi-hour manual process to a sub-30-second automated
alert --- directly addressing the failure mode that caused the April~17
incident. My Dynatrace parity analysis demonstrates that the AWS-native
stack achieves 60\% functional parity at approximately \$77/month, rising to
80\% with Application Signals, versus Dynatrace's enterprise licensing cost.

% ────────────────────────────────────────────────────────────
\section{Report Structure}
% ────────────────────────────────────────────────────────────

The remainder of this report is organized as follows.
Chapter~\ref{chap:tools} surveys the AWS services, development tools, and
AI-native tooling used throughout the internship.
Chapter~\ref{chap:project} presents the full project context, team structure,
Meridian platform, and JIRA scope.
Chapter~\ref{chap:learning_goals} states the technical and non-technical
learning goals established at the outset of the engagement.
Chapter~\ref{chap:methodologies} documents the SRE methodology applied:
golden signals, SLOs, gap analysis, and requirements.
Chapter~\ref{chap:implementation} covers the phased technical delivery,
system architecture, and design decisions.
Chapter~\ref{chap:weekly_breakdown} provides a phase-by-phase delivery
narrative organized across the 24-week internship.
Chapter~\ref{chap:conclusions} evaluates outcomes against learning goals and
provides a personal reflection.
Chapter~\ref{chap:achievements} presents the key deliverables and business
value realized.
Chapter~\ref{chap:challenges} documents the six engineering challenges
encountered and their mitigations.
Chapter~\ref{chap:future_work} outlines the planned extensions, with emphasis
on the auto-healing agent (IISS-488).
